\vspace{-5pt}

\section{Scope and Problem Statement}
\label{sec:scope}
\vspace{-3pt}


\noindent\textbf{Scope.}  
This paper focuses specifically on the problem of \emph{
    developer-style prompts generating malicious code snippets.}
The scope of this paper is limited to the following:  

\begin{itemize}[nosep,leftmargin=*]
    \item We only consider malicious code generated directly by LLMs. 
    We do not consider malicious content introduced by  
    external tools such as search engines.  
    \item We restrict attention to developer-style prompts that 
    could be asked in normal development tasks. 
    We do not consider adversarial prompting, jailbreaking, 
    prompt injection, and all other active inference-time attack techniques. 
\end{itemize}

While external tools can introduce contaminated content, this represents a separate
attack vector that has been explored in prior research, such as search engine
optimization. Furthermore, the presence of external contamination would only make the
security issues of LLMs worse. Although adversarial prompting and jailbreaking are
important methods used to actively exploit or
bypass an LLM's safety features at the moment a prompt is submitted, 
they form
a different threat model and they are much less 
likely to be used by regular developers, so we exclude them from our scope.

% \textbf{Scope.} The scope of this paper is limited to the following:
% \begin{itemize}
% \item We only consider the malicious code generated by LLMs themselves, 
% we do not consider the malicious content introduced in search engines, 
% or any other tools used by LLMs.~\footnote{Studying how to detect and prevent 
% malicious contents in search engine is beyond the scope of this paper and has been
% researched extensively. In fact, poisoned data introduced by other tools will 
% only make this security problem worse.}

% \item We do not consider adversarial prompting, jailbreaking, and prompt injection 
% techniques. We primarily investigate the malicious behavior of LLMs 
% when facing normal development tasks.  
% \end{itemize}

% Our focus on URLs in programming tasks is motivated by the example 
% presented in Section~\ref{sec:example_problem_statement}. 
% Malicious URLs in programming 
% contexts pose direct and immediate threats, potentially causing catastrophic 
% consequences including private information leakage, or cryptocurrency asset theft. 

% Our exclusion of adversarial prompting captures the majority of real-world 
% LLM usage where vulnerabilities have maximum impact potential. Because
% we focus on the code generated by LLMs in response to normal development tasks. 
% Because that's when developers will trust the generated code the most, 
% potentially leading to the execution of malicious code.



\noindent\textbf{Problem Statement.}  
Let $\mathcal{M}$ denote a large language model which takes prompts as input and 
generates code snippets as output, and let $\mathcal{O}$ denote an 
oracle function that determines 
whether a code snippet is malicious:  
$\mathcal{O}: \text{code} \rightarrow \{\text{benign}, \text{malicious}\}$.

We further assume the existence of an oracle $\mathcal{P}$ that classifies user prompts as either 
``developer-style'' (benign developer-style requests) or 
``non-developer-style'' (crafted to exploit model vulnerabilities
or phrased in ways unlikely to occur in normal software development practice):  
$\mathcal{P}: \text{prompt} \rightarrow \{\text{developer-style}, \text{non-developer-style}\}$.  

We define $\mathcal{S}$ as the set of prompt-code pairs where a developer-style prompt 
elicits a malicious code snippet from the model:
$\mathcal{S} = \{\, (p, c) \mid \mathcal{P}(p) = \text{developer-style}, \; 
c = \mathcal{M}(p), \; \mathcal{O}(c) = \text{malicious} \,\}$  

\noindent\textbf{Objective.}
Given $\mathcal{M}$, $\mathcal{O}$, and $\mathcal{P}$, our objective is to develop a framework 
to automatically discover and systematically expand the set $\mathcal{S}$.

In this paper, $\mathcal{O}$ is instantiated as an oracle that flags a code snippet as malicious 
if it contains at least one malicious URL.  
Human annotators serve as the oracle $\mathcal{P}$ to validate whether
a prompt is indeed a benign developer-style request.